\chapter{Onde aprender mais sobre Quarkus}
\label{cap:onde-aprender}

\proximasaulas

A documentação do Quarkus muda a cada versão, e os livros envelhecem
mais rápido que os de Java. A ordem sugerida: guias oficiais primeiro
(sempre na versão que você usa), livros para a visão de conjunto e o
Quarkus Insights para acompanhar o que está mudando.

\section{Guias oficiais}

Os guias citados nesta Parte, agrupados por assunto:
\begin{description}
\item[Build time] \citetitle{doc:quarkus-writing-extensions},
  \citetitle{doc:quarkus-cdi-reference},
  \citetitle{doc:quarkus-class-loading}.
\item[Arquitetura] \citetitle{doc:quarkus-cdi},
  \citetitle{doc:quarkus-rest}, \citetitle{doc:quarkus-panache},
  \citetitle{doc:quarkus-config-mappings}.
\item[Testes e operação] \citetitle{doc:quarkus-testing},
  \citetitle{doc:quarkus-dev-services},
  \citetitle{doc:quarkus-fault-tolerance},
  \citetitle{doc:quarkus-health}, \citetitle{doc:quarkus-micrometer},
  \citetitle{doc:quarkus-opentelemetry}.
\item[Decisões] \citetitle{doc:quarkus-building-native},
  \citetitle{doc:quarkus-native-tips},
  \citetitle{doc:quarkus-performance-measure},
  \citetitle{doc:quarkus-virtual-threads},
  \citetitle{doc:quarkus-reactive-architecture}.
\end{description}

Para quem quer entender o porquê do projeto, a página
\citetitle{doc:quarkus-container-first} resume a filosofia em poucas
linhas.

\section{Livros}

\small
\begin{longtable}{@{}p{4.6cm} p{3cm} p{4.6cm}@{}}
\toprule
Livro & Autor(es), ano & Por que ler \\
\midrule \endhead
\citetitle{livro:stefanko2025} & \citeauthor{livro:stefanko2025},
  \citeyear{livro:stefanko2025} & O mais recente da lista; visão de
  conjunto do desenvolvimento com Quarkus \\
\citetitle{livro:clingan2021} & \citeauthor{livro:clingan2021},
  \citeyear{livro:clingan2021} & Microsserviços com Quarkus e as APIs do
  MicroProfile em Kubernetes \\
\bottomrule
\end{longtable}
\normalsize

Os dois são da Manning, em inglês; não encontramos edição em português.
Confira a versão do Quarkus usada em cada um antes de copiar
configuração: o livro de 2021 é anterior às virtual threads finais do
JDK~21 (\citetitle{doc:jep444}, 2023), e a página da editora do livro de 2025 não
informa a versão coberta.

\section{Vídeos}

O \citetitle{video:quarkus-insights} é um podcast em vídeo semanal,
apresentado por desenvolvedores do Quarkus, sobre o que está mudando no
projeto e no ecossistema. Os episódios ficam no canal do Quarkus no
YouTube, linkado na página.

\section{Desafio prático}

\desafio{Escolha um guia oficial desta lista que você ainda não leu.
Leia na versão do Quarkus do seu projeto e escreva, em cinco linhas, uma
coisa que você faria diferente no seu serviço a partir dele, com o
trecho do guia que justifica.}
